\documentclass[14pt,dvipdfmx]{jsarticle}
\input{lecture-preamble.tex}

% 指数関数のグラフ用の方眼と座標軸（tikzpicture の中で使う）
% 使い方: \expaxes{y の最大値}
\newcommand{\expaxes}[1]{%
  \foreach \gx in {-3.5,-3,...,3.5} \draw[gray!15] (\gx,-1.5) -- (\gx,#1.5);
  \foreach \gy in {-1.5,-1,...,#1.5} \draw[gray!15] (-3.5,\gy) -- (3.5,\gy);
  \foreach \gx in {-3,...,3} \draw[gray!65] (\gx,-1.5) -- (\gx,#1.5);
  \foreach \gy in {-1,...,#1} \draw[gray!65] (-3.5,\gy) -- (3.5,\gy);
  \draw[-{Stealth}] (-3.7,0) -- (3.8,0) node[right]{$x$};
  \draw[-{Stealth}] (0,-1.7) -- (0,#1.8) node[above]{$y$};
  \node[below left, inner sep=1pt, font=\scriptsize] at (0,0) {O};
  \foreach \i in {-3,-2,-1,1,2,3} \node[below, font=\scriptsize] at (\i,0) {$\i$};
  \foreach \j in {1,...,#1} \node[left, font=\scriptsize] at (0,\j) {$\j$};
}

% 対応表のセル（幅をそろえる）: \tcell{値}，答えのセル: \tans{答え}
\newcommand{\tcell}[1]{\makebox[3.1zw]{$#1$}}
\newcommand{\tans}[1]{\makebox[3.1zw]{\ifshowanswer\answerstyle{$#1$}\fi}}

\begin{document}

\lectureheader{数学Ⅱ}{5}{指数関数と対数関数}{1}{指数関数}{2-A}{指数関数 $y=a^x$ のグラフ}
  {156,157}{指数関数のグラフを描けるようになろう}

% 同じ節の2枚目以降は，前回までの番号を引き継ぐ
% \setcounter{definition}{1}

% ============ 1ページ目（表・左） ============
\heading{〜前時の復習〜}
\begin{theorembox}{指数法則}
  $a > 0$，$b > 0$ とし，$r$，$s$ は有理数とする．
  \begin{tasks}[label=\arabic*, label-format=\bfseries, label-width=1zw, item-indent=1.5zw, after-item-skip=0.5em](2)
    \task $a^r \times a^s = \fitblankbf{a^{r+s}}$
    \task $\frac{a^r}{a^s} = \fitblankbf{a^{r-s}}$
    \task $(a^r)^s = \fitblankbf{a^{rs}}$
    \task $(ab)^r = \fitblankbf{a^r b^r}$
    \task $a^{-r} = \fitblankbf{\frac{1}{a^r}}$
  \end{tasks}
\end{theorembox}

\heading{復習問題}
\textbf{1}\quad 次の計算をせよ。
\begin{tasks}[label=(\arabic*), label-width=1.5zw, item-indent=2zw, after-item-skip=0.8em](2)
  \task $3^0 = \answermath{1}$
  \task $5^{-1} = \answermath{\frac{1}{5^1} = \frac{1}{5}}$
  \task* $2^2 \times 2^{-3} = \answermath{2^{2+(-3)} = 2^{-1} = \frac{1}{2}}$
\end{tasks}

\textbf{2}\quad 次の計算をせよ。
\begin{tasks}[label=(\arabic*), label-width=1.5zw, item-indent=2zw, after-item-skip=1.5em](1)
  \task $9^{\frac{1}{2}} = \answermath{(3^2)^{\frac{1}{2}} = 3^{2 \times \frac{1}{2}} = 3}$
  \task $125^{\frac{2}{3}} = \answermath{(5^3)^{\frac{2}{3}} = 5^{3 \times \frac{2}{3}} = 5^2 = 25}$
  \task $2^{-\frac{1}{3}} \times 2^{\frac{5}{6}}
    = \answermath{2^{-\frac{1}{3}+\frac{5}{6}} = 2^{-\frac{2}{6}+\frac{5}{6}} = 2^{\frac{3}{6}} = 2^{\frac{1}{2}} = \sqrt{2}}$
\end{tasks}

\vfill
\nextpage

% ============ 2ページ目（表・右） ============
\begin{definitionbox}{指数関数}
  $a$ を $1$ と異なる正の定数とするとき，$y=a^x$ は $x$ の関数である．
  この関数を，$a$ を\fitblankbf{\text{底}}とする $x$ の\fitblankbf{\text{指数関数}}という．
\end{definitionbox}
\vspace{0.5em}

\heading{指数関数 $y=2^x$ のグラフを調べてみよう！}
$x$ が $-2$，$-0.5$，$1.5$ のときの $2^x$ の値を求めよう．（ただし，$\sqrt{2} \fallingdotseq 1.414$ とする．）
\par\vspace{0.5em}
\hspace*{3zw}$\begin{alignedat}{2}
  & 2^{-2}   &&= \answermath{\frac{1}{2^2} = \frac{1}{4} = 0.25} \\[1.2em]
  & 2^{-0.5} &&= \answermath{2^{-\frac{1}{2}} = \frac{1}{\sqrt{2}} = \frac{\sqrt{2}}{2} \fallingdotseq 0.71} \\[1.2em]
  & 2^{1.5}  &&= \answermath{2^{\frac{3}{2}} = 2^{1+\frac{1}{2}} = 2\sqrt{2} \fallingdotseq 2.83}
\end{alignedat}$
\par\vspace{2em}

$y=2^x$ における $x$ と $y$ の対応表の空らんを埋めよう．
\begin{center}
  \renewcommand{\arraystretch}{1.5}\setlength{\tabcolsep}{3pt}
  \setlength{\arrayrulewidth}{0.8pt}\setlength{\doublerulesep}{1.5pt}
  \begin{tabular}{|c||c|c|c|c|c|c|c|c|c|}
    \hline
    $x$ & \tcell{-2} & \tcell{-1.5} & \tcell{-1} & \tcell{-0.5} & \tcell{0}
        & \tcell{0.5} & \tcell{1} & \tcell{1.5} & \tcell{2} \\ \hline
    $y$ & \tcell{0.25} & \tans{0.35} & \tans{0.5} & \tcell{0.71} & \tans{1}
        & \tans{1.41} & \tcell{2} & \tcell{2.83} & \tcell{4} \\ \hline
  \end{tabular}
\end{center}
\vspace{1em}

\sidebyside[0.45]{%
  座標平面上で，上の表の $x$，$y$ の値の組を座標にもつ点 $(x, y)$ をとると，
  右の図のような曲線上にある．\par\vspace{0.5em}
  この曲線が指数関数 $y=2^x$ のグラフである．
}{%
  \begin{tikzpicture}[x=0.65cm, y=0.65cm]
    \expaxes{6}
    \ifshowanswer
      \draw[blue, thick, domain=-3.5:2.7, samples=60] plot (\x, {2^(\x)});
      \foreach \p in {-2,-1.5,...,2} \fill[blue] (\p, {2^(\p)}) circle (2pt);
      \node[blue, above left] at (-1.6,0.4) {$y=2^x$};
    \fi
  \end{tikzpicture}
}

\vfill
\nextpage

% ============ 3ページ目（裏・左） ============
\heading{指数関数 $y=\left(\frac{1}{2}\right)^x$ のグラフを調べてみよう！}
\par\vspace{0.5em}
\hspace*{1zw}$\displaystyle \left(\frac{1}{2}\right)^x = \answermath{\frac{1}{2^x} = 2^{-x}}$ であるから，\par
\hspace*{1zw}$\displaystyle y=\left(\frac{1}{2}\right)^x$ は $y=\fitblankbf{2^{-x}}$ と同じである．
\par\vspace{1em}
{%
  \setlength{\lineskiplimit}{6pt}\setlength{\lineskip}{8pt}%
  したがって，$y=\left(\frac{1}{2}\right)^x$ のグラフは，
  $y=2^x$ のグラフと \fitblankbf{y \text{軸}}に関して対称であり，下の図のようになる．\par
}

\begin{center}
  \begin{tikzpicture}[x=0.7cm, y=0.7cm]
    \expaxes{6}
    \draw[gray, thick, dashed, domain=-3.5:2.7, samples=60] plot (\x, {2^(\x)});
    \node[gray, right] at (2.7,6) {$y=2^x$};
    \ifshowanswer
      \draw[blue, thick, domain=-2.7:3.5, samples=60] plot (\x, {2^(-\x)});
      \node[blue, left] at (-2.7,6) {$y=\left(\frac{1}{2}\right)^x$};
    \fi
  \end{tikzpicture}
\end{center}

\vfill
\textbf{補足}\quad 一般に，関数 $y=f(x)$ のグラフと $y=f(-x)$ のグラフは，
$y$ 軸に関して対称である．
\par\vspace{0.5em}
\begin{techniquebox}{グラフの対称移動}
  関数 $y=f(x)$ のグラフに対して，
  \begin{itemize}[leftmargin=1.5zw, itemsep=0.2em]
    \item $y=f(-x)$ のグラフは\fitblankbf{y \text{軸}}に関して対称
    \item $y=-f(x)$ のグラフは\fitblankbf{x \text{軸}}に関して対称
    \item $y=-f(-x)$ のグラフは\fitblankbf{\text{原点}}に関して対称
  \end{itemize}
\end{techniquebox}

\nextpage

% ============ 4ページ目（裏・右） ============
\heading{指数関数 $y=a^x$ のグラフ}
\par\noindent
\begin{minipage}[t]{0.48\linewidth}
  \textbf{(i)\enspace\boldmath$a>1$ のとき}\par
  \centering
  \begin{tikzpicture}[x=0.5cm, y=0.45cm]
    \draw[-{Stealth}] (-3,0) -- (2.6,0) node[right]{$x$};
    \draw[-{Stealth}] (0,-0.6) -- (0,4.4) node[above]{$y$};
    \node[below left, font=\footnotesize] at (0,0) {O};
    \ifshowanswer
      \begin{scope}[blue]
        \draw[thick, domain=-3:1.75, samples=60] plot (\x, {2.2^(\x)});
        \draw[dashed] (1,0) node[below, font=\footnotesize]{$1$} -- (1,2.2) -- (0,2.2) node[left, font=\footnotesize]{$a$};
        \node[left, font=\footnotesize] at (0,1) {$1$};
      \end{scope}
    \fi
  \end{tikzpicture}
\end{minipage}\hfill
\begin{minipage}[t]{0.48\linewidth}
  \textbf{(ii)\enspace\boldmath$0<a<1$ のとき}\par
  \centering
  \begin{tikzpicture}[x=0.5cm, y=0.45cm]
    \draw[-{Stealth}] (-2.6,0) -- (3,0) node[right]{$x$};
    \draw[-{Stealth}] (0,-0.6) -- (0,4.4) node[above]{$y$};
    \node[below left, font=\footnotesize] at (0,0) {O};
    \ifshowanswer
      \begin{scope}[blue]
        \draw[thick, domain=-1.75:3, samples=60] plot (\x, {0.45^(\x)});
        \draw[dashed] (1,0) node[below, font=\footnotesize]{$1$} -- (1,0.45) -- (0,0.45) node[left, font=\footnotesize]{$a$};
        \node[left, font=\footnotesize] at (0,1) {$1$};
      \end{scope}
    \fi
  \end{tikzpicture}
\end{minipage}
\par\vspace{0.3em}

\begin{definitionbox}{漸近線}
  曲線が限りなく近づいていく直線を，その曲線の\fitblankbf{\text{漸近線}}という．
\end{definitionbox}

\begin{theorembox}{指数関数 $y=a^x$ のグラフの特徴}
  \begin{enumerate}[label=\textbf{\arabic*}, leftmargin=1.5zw, itemsep=0.2em]
    \item 2点 $(0,\ \fitblankbf{1})$，$(1,\ \fitblankbf{a})$ を通る．
    \item $\fitblankbf{x \text{軸}}$ を漸近線としてもつ．
    \item $a>1$ のとき\fitblankbf{\text{右上がり}}の曲線，\\
          $0<a<1$ のとき\fitblankbf{\text{右下がり}}の曲線である．
  \end{enumerate}
\end{theorembox}
\vspace{0.2em}

\begin{summarybox}
  \begin{enumerate}[label=\textbf{\arabic*}, leftmargin=1.5zw, itemsep=0.2em]
    \item $a>0$，$a \neq 1$ のとき，$y=a^x$ を $a$ を\fitblankbf{\text{底}}とする\\ $x$ の\fitblankbf{\text{指数関数}}という．
    \item $y=\left(\frac{1}{a}\right)^x$ のグラフは，$y=a^x$ のグラフと\fitblankbf{y \text{軸}}に関して対称である．
    \item $y=a^x$ のグラフは点 $(0,\ \fitblankbf{1})$ を通り，\fitblankbf{x \text{軸}}を漸近線とする．\\
          $a>1$ のとき\fitblankbf{\text{右上がり}}の曲線，\\
          $0<a<1$ のとき\fitblankbf{\text{右下がり}}の曲線である．
  \end{enumerate}
\end{summarybox}

\heading{確認問題}
教科書p.157 練習9

\vfill
\nexttime{指数関数の特徴を踏まえて，指数を含む方程式，不等式について考えてみよう！}

\end{document}
